\chapter{背景与动机}

\section{MNN 执行机制与移动端运行环境}

\subsection{MNN 在线程池与 CPU 后端中的执行路径}

MNN 的 CPU 推理执行路径建立在后端抽象、线程池和算子执行器之上。对本文相关的优化而言，最关键的入口位于 CPUBackend、线程池封装以及并行宏定义层。算子在执行时，并不是直接感知“当前是 Prefill 还是 Decode”，而是通过后端配置、线程数和任务划分计划间接决定并发行为。换句话说，若要让不同阶段采用不同调度策略，就必须把阶段信息、线程参数和 affinity 参数准确传到 CPU 后端。

从本文代码实现看，阶段相关状态主要由 AutoTuner 维护。运行时通过设置当前阶段、当前线程数和当前 affinity 掩码，让 CPUBackend 在任务划分时获得正确的上下文。随后，并行宏将具体策略展开为线程池任务：对于 Prefill，系统初始化 work stealing 状态，再由各线程循环获取本地块或窃取块；对于 Decode，系统则使用更轻量的动态获取逻辑。这样的组织方式有两个优点。第一，不需要彻底改写 MNN 原有的后端结构。第二，调度策略可以在阶段切换时快速切换，而不必为每个算子单独写一套外层控制逻辑。

这一设计与许多移动端推理优化工作形成呼应。已有研究表明，通用推理框架如果忽略底层处理器异构性，就很容易出现线程扩展性差和资源利用不均的问题\cite{wang2021asymo,xu2021coschedule,kang2020embedded}。因此，把调度信息接入框架运行时，而不是停留在应用层脚本，是系统优化真正生效的前提。

\subsection{Android/Linux 中的 affinity 与线程迁移}

移动端推理系统最终运行在 Android/Linux 调度器之上。应用层能控制的范围有限，但并非什么都做不了。即使没有 root 权限，用户态程序仍然可以通过标准 affinity 接口约束线程运行的 CPU 集合。对本文而言，这一点非常重要，因为它保证了“分阶段绑核”可以在本地环境下落地，而不依赖修改系统调度器或内核参数。

如果不显式设定 affinity，线程会由操作系统按全局负载自动分配。这样的默认策略对系统整体吞吐是友好的，但未必适合单个深度学习任务。推理线程在不同核心之间迁移后，L1/L2 缓存内容往往无法复用，之前建立的局部性会被打断，线程还可能落到性能较低的核心上继续执行。对于 Prefill 这类长任务阶段，迁移造成的性能损失会持续累积；对于 Decode 这类短步迭代阶段，迁移会放大每一步的延迟抖动。

已有面向异构移动 CPU 的研究已经指出，线程分工与核心绑定必须共同考虑\cite{wang2021asymo,abdelgawad2022bertperf}。仅仅提高线程数并不能自动换来更好性能；若线程分布与核心能力不匹配，反而会加剧争用和拖尾。因此，本文把 affinity 视为和线程数同等重要的一级参数，而不是可有可无的附属配置。

\subsection{无 root 条件下的 ADB 与 Perfetto 观测能力}

本文实验明确限定在无 root 条件下进行。这一约束会影响系统可观测性，但并不会使研究无法开展。首先，通过 ADB 可以批量执行推理测试脚本、收集日志并控制实验流程。其次，Perfetto 在普通用户权限下仍然可以提供线程调度、CPU 运行时间片和部分系统事件轨迹。这些信息已经足以支持本文对“线程迁移是否减少”“高负载时是否出现尾部阻塞”“work stealing 是否缓解长尾”的分析。

需要指出的是，无 root 条件下无法方便地读取全部功耗和硬件计数器信息，因此本文不以能耗为核心评价目标。这不是实验缺失，而是设计边界的主动收缩：本文聚焦可在普通开发环境下复现的性能优化路径，而不把论证建立在需要特权访问的数据之上。与之对应，实验指标将集中在时延、吞吐、稳定性和调度行为这四个方面。

\section{Prefill 与 Decode 的阶段差异}

\subsection{Prefill 阶段的长任务与负载不均}

在大语言模型推理中，Prefill 负责处理输入提示词并建立上下文状态。该阶段通常面对较长序列，单次执行包含更多算子调用和更大的张量规模。由于总工作量大，系统更关注整体吞吐和负载均衡。如果任务划分不均，尾部线程会显著拖慢整个阶段完成时间。

这类阶段与经典 work stealing 场景有天然相似性：任务总体较大，局部不均衡会在运行过程中不断暴露，而适度的任务窃取可以把已空闲线程重新利用起来\cite{blumofe1999workstealing,acar2000locality}。但与通用并行程序不同，移动端推理还受到异构核心能力差异和缓存局部性的制约。因此，单纯随机 stealing 并不一定最优，它必须建立在较合理的初始划分之上。

从本文代码实现看，Prefill 阶段先按核心能力做加权初始划分，再由各线程优先执行自己的本地队列，只有本地工作耗尽后才触发窃取。这种“先静态、后 stealing”的思路，本质上是在低开销和高适应性之间求平衡。若一开始就完全依赖动态调度，线程之间的竞争和同步成本会偏高；若完全依赖静态划分，又无法处理高负载和性能比失准带来的长尾。

\subsection{Decode 阶段的短步迭代与同步开销}

Decode 阶段的特点与 Prefill 有明显差异。此时模型按 token 逐步生成，单步任务规模更小，调用频率更高。系统往往不再追求一次性处理大量工作，而是追求每一步都尽量短、尽量稳。换言之，Decode 的核心问题不是“如何把总任务充分分散”，而是“如何避免调度本身成为额外负担”。

这一点和服务系统领域对 Prefill/Decode 差异的观察是一致的。Orca、PagedAttention、DistServe 和 Sarathi 等工作都指出，预填充和生成阶段的资源需求不同，若把它们视为完全同质的负载，系统性能往往会受限\cite{yu2022orca,kwon2023pagedattention,zhong2024distserve,agrawal2023sarathi}。虽然这些工作主要面向服务端，但其关于阶段差异的结论同样适用于端侧。

因此，本文没有把 Decode 也强行设计成和 Prefill 完全相同的 work stealing 流程，而是采用更轻量的动态任务获取方式。这样做的原因很直接：Decode 的单步任务短，若每一步都进行复杂的窃取判定和多队列竞争，调度开销本身就可能抵消收益。阶段感知调度的含义，正在于允许系统根据阶段特征选择不同复杂度的机制。

\subsection{阶段独立线程数和绑核的必要性}

一旦承认 Prefill 和 Decode 不是同一种负载，就必须进一步承认两者不应共用一套参数。Prefill 更看重大核和总吞吐，Decode 更看重执行稳定性和每步开销。前者通常可以接受更强的并行和更积极的 stealing，后者则更适合较轻的调度逻辑和更谨慎的线程扩展。

本文前期实验表明，两阶段性能曲线的最优线程数并不重合。相应地，两阶段的最优 CPU 集合也可能不同。对计算密集型阶段，可以优先选取高性能核心；对更偏访存敏感、对尾延迟更敏感的阶段，则可能需要在性能和稳定性之间权衡。把这种差异显式建模为“阶段独立线程数”和“阶段独立 affinity”，是本文系统设计的基本前提。

\section{现有任务划分与调度方法分析}

\subsection{静态划分、guided scheduling 与 work stealing}

静态划分的优点是简单、开销低、可预测性强。在线程能力相近、任务规模均匀的场景中，静态划分往往已经足够好。但在异构多核和运行时扰动明显的环境下，静态划分的弱点也很突出：一旦初始边界不合理，后续几乎没有自我修正能力。

guided scheduling 在此基础上引入了逐步缩小任务块的思想。其核心目的是在保留较低调度开销的同时，减轻尾部不均衡\cite{polychronopoulos1987guided}。这一思想对本文有启发意义，但 Prefill 阶段面临的不是抽象对称处理器，而是异构多核和高负载扰动。相比之下，work stealing 更适合处理“初始分工基本合理，但运行过程中仍可能出现局部失衡”的场景\cite{blumofe1999workstealing,frigo1998cilk5}。

work stealing 的经典设计是每个工作线程维护本地双端队列，线程优先从本地队列取任务；当本地队列为空时，再去窃取其他线程队列头部或尾部的任务。Chase-Lev deque 和针对弱内存模型的后续工作，进一步给出了这种结构在无锁并发环境中的工程实现方式\cite{chase2005deque,le2013weakmemory}。这为本文在 ARM 平台上实现无锁任务窃取提供了直接参考。

\subsection{异构多核上的性能比建模与加权划分}

异构多核任务划分不能只回答“划多少”，还要回答“按什么比例划”。移动端推理相关研究已经反复说明，异构处理器上最大的误区之一，就是把不同核心当作同构资源使用\cite{wang2021asymo,liu2024pitfalls}。若忽略核心能力差异，线程扩展性会迅速恶化；若只使用固定经验值，又可能在不同模型和不同机型上失准。

AsyMo 这类研究给出的启示是：异构调度应尽可能建立在目标负载和目标硬件的可测特征上，而不是完全依赖通用经验参数\cite{wang2021asymo}。BERTPerf 则从另一个角度说明，big.LITTLE 上的推理延迟与配置参数之间存在复杂关系，必须通过建模或实测加以刻画\cite{abdelgawad2022bertperf}。这些工作与本文的联系在于，它们都强调“真实能力测量”比“统一经验参数”更可靠。

因此，本文在静态划分之前增加了核心性能比标定步骤。它并不追求复杂的全局最优模型，而是服务于一个更直接的目标：让初始任务边界尽量符合目标机型的真实速度关系。后续的 work stealing 则负责处理仍然存在的残余失衡。

\subsection{移动端推理框架中的已有不足}

从现有框架和相关工作看，移动端推理优化主要集中在算子实现、量化、异构协同和调度建模几个方向\cite{jia2022codl,xu2021coschedule,niu2024smartmem}。这些工作提供了大量有价值的经验，但若回到“无 root、本地可运行、直接落到 MNN”的目标上，仍然能看到三个缺口。

第一，很多优化建立在通用异构协同或特定硬件后端上，未必直接适配只使用 CPU 的端侧部署路径。第二，许多研究强调离线建模和全局计划，但对高负载干扰下的运行时纠偏讨论不够充分。第三，面向大模型推理时，Prefill 和 Decode 的阶段差异常被服务系统研究关注，却较少被直接映射到移动端 CPU 后端调度中。

本文的工作正是围绕这三个缺口展开：以 MNN 为目标框架，以 CPU 为目标后端，以普通开发环境为目标约束，在不重构整体框架的前提下实现分阶段绑核、性能比加权划分和 Prefill work stealing，并通过真实手机实验验证其收益。

\section{问题归纳与设计目标}

\subsection{设计目标}

结合前述分析，本文将优化目标概括为三个方面。其一，在无 root 条件下提升端侧推理性能，缩短 Prefill 和 Decode 的执行时间。其二，在背景负载存在时维持较好的稳定性，减少长尾和明显抖动。其三，在保持实现简洁的前提下，将优化逻辑嵌入现有 MNN 执行路径，使方案具备可复现性和工程可用性。

\subsection{约束条件}

本文的设计并不试图控制整个系统调度器，也不依赖特权接口读取功耗或修改 governor。优化必须在普通用户权限下完成，能够通过应用层或用户态运行时直接生效。与此同时，本文也不对 MNN 做大规模结构重写，而是尽量利用现有线程池、后端接口和并行宏定义完成最小侵入式改造。

\subsection{评价指标}

基于上述目标和约束，本文的评价指标主要包括：Prefill 阶段吞吐、Decode 阶段吞吐、首 token 时延、端到端时延、不同重复实验下的稳定性统计，以及基于日志和 Perfetto 的线程调度行为指标。对于 work stealing，本研究还将关注窃取尝试次数、窃取成功次数和被窃取任务量等统计量，以便把性能变化与调度行为直接对应起来。
